\documentclass[11pt,a4paper]{article}
\usepackage[T1]{fontenc}
\usepackage[utf8]{inputenc}
\usepackage{lmodern}
\usepackage[a4paper,left=2.3cm,right=2.1cm,top=2.2cm,bottom=2.1cm]{geometry}
\usepackage{amsmath,booktabs,longtable,tabularx,array,multirow}
\usepackage{graphicx,xcolor,microtype,enumitem,ragged2e,float,pdflscape}
\usepackage{caption,fancyhdr,titlesec}
\usepackage[hidelinks]{hyperref}
\renewcommand{\contentsname}{Sumário}
\renewcommand{\refname}{Referências}
\renewcommand{\tablename}{Tabela}

\definecolor{azul}{HTML}{2A78D6}
\definecolor{laranja}{HTML}{EB6834}
\definecolor{tinta}{HTML}{14181F}
\definecolor{tinta2}{HTML}{525B66}
\definecolor{linha}{HTML}{DFE3E8}
\definecolor{fundo}{HTML}{F4F5F7}
\captionsetup{font=small,labelfont={bf,color=tinta},textfont={color=tinta2},
  justification=justified,singlelinecheck=false,skip=5pt}
\titlespacing*{\section}{0pt}{12pt plus 2pt minus 2pt}{5pt}
\titlespacing*{\subsection}{0pt}{9pt plus 2pt minus 2pt}{3pt}
\setlength{\parskip}{2pt plus 1pt}
\setlist{topsep=3pt,itemsep=2pt,parsep=1pt}
\newcolumntype{Y}{>{\RaggedRight\arraybackslash}X}
\newcommand{\destaque}[2]{\begin{center}\begin{minipage}{.96\textwidth}
  \colorbox{fundo}{\begin{minipage}{\dimexpr\linewidth-2\fboxsep}
  \textbf{\color{azul}#1}\par\small #2\end{minipage}}\end{minipage}\end{center}}

\newcommand{\numerogrupo}{[substituir]}

\pagestyle{fancy}\fancyhf{}
\fancyhead[L]{\small\color{tinta2}Banco de Dados --- DAS e Arquiteturas RAID}
\fancyfoot[C]{\small\thepage}
\renewcommand{\headrulewidth}{.4pt}
\hypersetup{pdftitle={Arquiteturas de Armazenamento Fisico para SBD: DAS e RAID},
  pdfauthor={Grupo de Banco de Dados --- UFRJ}}

\begin{document}
\thispagestyle{empty}
\begin{center}
{\large\scshape Universidade Federal do Rio de Janeiro}\\[2pt]
{\normalsize Instituto de Matemática --- Departamento de Ciência da Computação}\\[2pt]
{\normalsize Disciplina: Construção de Banco de Dados}\\[2pt]
{\normalsize Professor: Milton Ramirez}\\[17mm]
{\color{azul}\rule{\textwidth}{1.2pt}}\\[6mm]
{\LARGE\bfseries Cloud Computing e Cloud Storage}\\[4mm]
{\Large Estudo de caso: Amazon Web Services (AWS)}\\[4mm]
{\large Arquitetura, serviços, segurança e análise de 3,5 PB}\\[6mm]
{\color{azul}\rule{\textwidth}{1.2pt}}\\[8mm]
{\large\bfseries Grupo \numerogrupo}\\[6mm]
\begin{tabular}{@{}ll@{}}\toprule
\textbf{Nome completo} & \textbf{DRE}\\\midrule
Bernardo Brandão Pozzato Carvalho Costa & 123289593\\
Enzo de Carvalho Sampaio & 123386206\\
Gabriel Schmitz Corrêa Rizawinsk & 123225573\\
Guilherme En Shih Hu & 123224674\\
Raphael Henrique da Silva Pereira & 123420783\\
Vivian Maria da Silva e Souza & 123205793\\\bottomrule
\end{tabular}\\[12mm]
{\normalsize Rio de Janeiro --- Setembro de 2026}
\end{center}
\clearpage

\tableofcontents
\clearpage

\section{Resumo}
Este trabalho apresenta um estudo sobre arquiteturas de armazenamento físico (storage) para Sistemas de Banco de Dados (SBD), com foco em Direct Attached Storage (DAS) e em arranjos RAID (Redundant Array of Independent Disks). São discutidas as principais interfaces de conexão para DAS --- SATA, eSATA, SAS, FireWire, USB (3 e 4) e HBA (Host Bus Adapter) ---, comparando velocidade, topologia e adequação a cargas de trabalho de banco de dados. Em seguida, aprofunda-se o estudo de RAID: os conceitos de paralelismo (desempenho) e redundância (disponibilidade e confiabilidade); as técnicas de data striping em nível de bit e de bloco; mirroring e shadowing; as métricas de confiabilidade MTTF, MTBF, MTTR e MTTDL; os níveis padrão de RAID (0 a 6, além dos arranjos aninhados 0+1 e 10); e níveis não padrão (RAID 1.5, RAID 7, RAID-DP, RAID-S, Matrix RAID, entre outros). Por fim, apresentam-se recomendações comparativas de uso e aplicação de cada nível conforme o perfil de carga típico de um SBD. O trabalho foi produzido com apoio de um assistente de Inteligência Artificial generativa (Claude, da Anthropic), conforme autorizado pelo enunciado da tarefa, e inclui um relatório Post-Mortem detalhando o uso dessa ferramenta.

\section{Introdução}
Todo Sistema de Banco de Dados depende, em última instância, de um subsistema de armazenamento físico capaz de gravar e recuperar dados de forma confiável, mesmo diante de falhas de hardware. Diferentemente da memória principal (volátil e de acesso rapidíssimo), os discos --- magnéticos (HDD) ou de estado sólido (SSD) --- são o meio de armazenamento não volátil onde residem as tabelas, índices e logs de um SGBD, e sua organização física tem impacto direto sobre desempenho, disponibilidade e durabilidade dos dados.

Neste trabalho, o grupo optou por estudar as arquiteturas de armazenamento do tipo DAS (Direct Attached Storage), em que os discos estão conectados diretamente ao servidor (ou a um host adapter dedicado), sem passar por uma rede de armazenamento compartilhada. Sobre essa base física, aprofunda-se o estudo das técnicas RAID, criadas justamente para combinar múltiplos discos físicos independentes em um único volume lógico, obtendo ganhos de desempenho (paralelismo), de confiabilidade (redundância) ou de ambos, a depender do nível escolhido.

O texto está organizado da seguinte forma: a Seção 3 situa o conceito de DAS frente a NAS e SAN; a Seção 4 descreve as interfaces de conexão SATA, eSATA, SAS, FireWire, USB e HBA; a Seção 5 introduz RAID e os conceitos de paralelismo e redundância; a Seção 6 detalha as técnicas de striping, mirroring e paridade; a Seção 7 apresenta as métricas de confiabilidade MTTF, MTBF, MTTR e MTTDL; a Seção 8 descreve comparativamente os níveis padrão de RAID (0 a 6, 0+1 e 10); a Seção 9 aborda níveis não padrão (RAID 1.5, 7, DP, S, Matrix, entre outros); a Seção 10 traz recomendações de uso para cenários típicos de SBD; a Seção 11 conclui o trabalho; e a Seção 12 apresenta o relatório Post-Mortem sobre o uso de IA generativa exigido pelo enunciado.

\section{Armazenamento Físico em SBD: DAS, NAS e SAN}
Direct Attached Storage (DAS) é a arquitetura de armazenamento em que um ou mais discos (ou um arranjo RAID de discos) estão conectados diretamente a um único servidor ou estação de trabalho, por meio de uma interface local --- um cabo SATA, SAS, USB, ou um HBA/controladora instalado no próprio barramento do servidor ---, sem intermediação de uma rede de armazenamento dedicada. É a forma mais simples e mais antiga de anexar storage a um SBD: os discos internos de um servidor de banco de dados, ou um gabinete externo de discos ligado por SAS a uma controladora RAID, são exemplos típicos de DAS.

Essa arquitetura contrasta com duas outras formas de organizar o armazenamento em rede, mencionadas aqui apenas para contextualizar a escolha do escopo deste trabalho:
\begin{itemize}[leftmargin=*]
    \item \textbf{NAS (Network Attached Storage):} dispositivo de armazenamento conectado à rede local (Ethernet/IP) e acessado por múltiplos hosts em nível de arquivo, por meio de protocolos como NFS ou SMB/CIFS --- o cliente enxerga arquivos e diretórios, não blocos brutos.
    \item \textbf{SAN (Storage Area Network):} rede dedicada de alto desempenho (tipicamente Fibre Channel ou iSCSI) que conecta múltiplos servidores a um pool compartilhado de armazenamento em nível de bloco, permitindo que vários hosts compartilhem os mesmos discos de forma centralizada e escalável.
\end{itemize}

Frente a essas duas alternativas, o DAS tem como principal vantagem a menor latência (não há overhead de rede entre o host e o disco) e a simplicidade de implantação e custo, o que o torna atrativo para servidores de banco de dados isolados ou de médio porte. Como desvantagem, o DAS tipicamente não permite que múltiplos servidores compartilhem o mesmo conjunto de discos com a mesma facilidade que uma SAN, e sua escalabilidade de capacidade fica limitada ao número de portas/baias do host ou da controladora. É sobre essa arquitetura --- DAS --- e sobre as interfaces e técnicas RAID usadas para torná-la ao mesmo tempo rápida e confiável que este trabalho se aprofunda.

\section{Interfaces de Conexão para DAS}
A conexão física entre o host e os discos de um arranjo DAS pode ser feita por diferentes interfaces, cada uma com características próprias de velocidade, topologia e adequação a cargas de trabalho de banco de dados. Esta seção descreve as interfaces indicadas no enunciado da tarefa: SATA, eSATA, SAS, FireWire, USB (3 e 4) e HBA.

\subsection{SATA (Serial ATA) e eSATA}
O SATA (Serial Advanced Technology Attachment) é a interface serial que substituiu o antigo PATA/IDE paralelo, tornando-se a interface padrão para discos de baixo custo em desktops e servidores de entrada. É uma interface ponto a ponto --- cada porta SATA conecta a controladora a um único disco --- que evoluiu em três gerações principais: SATA I (1,5 Gbps), SATA II (3 Gbps) e SATA III (6 Gbps, cerca de 550-600 MB/s efetivos), sendo essa última a revisão amplamente utilizada atualmente. O SATA suporta NCQ (Native Command Queuing), que permite reordenar até 32 comandos pendentes para otimizar o padrão de acesso ao disco, e é hot-swappable (permite troca do disco com o sistema em operação, quando suportado pelo chassi/controladora).

O eSATA (external SATA) é a variante externa do SATA: usa o mesmo protocolo e a mesma velocidade elétrica da interface interna, porém com um conector mais robusto e blindado, adequado para gabinetes externos. Sua principal limitação é não transportar energia (diferentemente do USB), exigindo um cabo de alimentação separado para o disco externo --- o que, somado à popularização do USB 3.x e do USB4 (com velocidades iguais ou superiores e fornecimento de energia integrado), fez o eSATA cair em disuso na maioria dos equipamentos modernos.

\subsection{SAS (Serial Attached SCSI)}
O SAS (Serial Attached SCSI) é a interface serial de classe empresarial, evoluída do antigo SCSI paralelo, mantendo seu conjunto de comandos (SCSI command set) sobre uma camada física serial. É a interface predominante em servidores de banco de dados e controladoras RAID corporativas, por três características que a diferenciam do SATA: dual-porting (cada disco SAS pode ser acessado por dois caminhos/controladoras independentes simultaneamente, permitindo multipath I/O e maior tolerância a falhas de controladora), suporte a múltiplos iniciadores por meio de expansores SAS (permitindo topologias com até 65.535 dispositivos endereçáveis) e uma fila de comandos mais profunda (até 254 comandos pendentes, contra cerca de 32 do NCQ do SATA). O SAS também é compatível com discos SATA (uma controladora SAS consegue endereçar discos SATA via STP --- SATA Tunneling Protocol), embora o inverso não seja possível.

As gerações do SAS dobraram de velocidade a cada revisão: SAS-1 (3 Gbps, 2004), SAS-2 (6 Gbps, 2009), SAS-3 (12 Gbps, 2013) e a geração atual, SAS-4 --- também chamada de 24G SAS ---, que opera a 22,5 Gbps por via (aproximadamente 2.400 MB/s), padronizada a partir de 2019. Dado seu foco em confiabilidade, redundância de caminho e desempenho sustentado sob carga contínua (certificação para operação 24/7), o SAS é a interface preferida para os discos que compõem arranjos RAID em servidores de SBD de produção.

\subsection{FireWire (IEEE 1394)}
O FireWire (padronizado como IEEE 1394 e também comercializado como i.LINK pela Sony) foi uma interface serial desenvolvida pela Apple nos anos 1990, historicamente usada para armazenamento externo e equipamentos de áudio e vídeo. Diferentemente do USB da época, o FireWire permitia comunicação peer-to-peer entre dispositivos (sem exigir que um computador atuasse como controlador central) e suportava encadeamento (daisy-chaining) de até 63 dispositivos. Suas duas revisões mais difundidas foram o FireWire 400 (400 Mbps) e o FireWire 800 (800 Mbps). Com a popularização do USB 3.x --- mais rápido, mais barato e universalmente suportado --- e, posteriormente, do Thunderbolt e do USB4, o FireWire caiu em desuso: hoje é considerado uma interface legada, raramente encontrada em equipamentos novos, e é mencionada aqui principalmente por seu valor histórico e por eventualmente ainda aparecer em DAS herdado de estações de trabalho e sistemas de captura de áudio/vídeo mais antigos.

\subsection{USB (USB 3 e USB4)}
O USB (Universal Serial Bus) é a interface mais universal para armazenamento externo, por sua ubiquidade, baixo custo e capacidade de fornecer energia ao dispositivo pelo próprio cabo. É, no entanto, uma interface host-cêntrica (hierárquica): diferentemente do SAS/SCSI, sempre exige um controlador host mediando a comunicação, o que a torna adequada principalmente para DAS de host único (backup, discos externos portáteis), e não para topologias multi-iniciador.

O USB 3.0 (também chamado USB 3.2 Gen 1), lançado em 2008/2011, introduziu o modo "SuperSpeed" de 5 Gbps; revisões seguintes elevaram a velocidade para 10 Gbps (USB 3.1 Gen2 / USB 3.2 Gen2) e 20 Gbps (USB 3.2 Gen2x2, usando dois canais). O USB4, lançado em 2019 e baseado no protocolo do Thunderbolt 3, unificou dados, vídeo (DisplayPort) e PCI Express em um único túnel, com velocidade de até 40 Gbps. A revisão mais recente, USB4 Versão 2.0 (publicada pelo USB-IF em outubro de 2022), praticamente dobrou esse limite: até 80 Gbps de forma simétrica (bidirecional) e, em um modo assimétrico opcional voltado a cargas como displays de altíssima resolução, até 120 Gbps em um sentido com 40 Gbps no sentido oposto --- usando uma nova camada física com sinalização PAM3. Essas velocidades aproximam o USB4 v2.0 do desempenho de interfaces internas de SSD, tornando viável, cada vez mais, um DAS externo de altíssimo desempenho conectado por uma simples porta USB-C.

\subsection{HBA (Host Bus Adapter)}
Diferentemente das interfaces anteriores, o HBA (Host Bus Adapter) não é, em si, um padrão físico de cabo, mas sim a placa controladora --- tipicamente um cartão PCI Express instalado no servidor --- que faz a ponte entre o barramento interno do host e a interface de armazenamento externa (SAS, SATA ou, em outros contextos, Fibre Channel), fornecendo as portas fisikcas de conexão aos discos e implementando o protocolo de comunicação com eles. É por meio do HBA que o restante das interfaces discutidas nesta seção efetivamente chega ao processador e à memória do servidor.

Para fins de RAID, a distinção mais importante é entre um HBA simples (também chamado de controladora "JBOD" ou passthrough, que apenas expõe cada disco individualmente ao sistema operacional) e uma controladora RAID de hardware, que agrega ao HBA um processador dedicado com engine de XOR, memória cache (frequentemente protegida por bateria ou capacitor com flash de backup --- BBU/FBWC) e firmware capaz de calcular paridade e reconstruir arranjos de forma autônoma, sem sobrecarregar a CPU do host. Quando não há uma controladora RAID de hardware, o RAID pode ainda assim ser implementado em software, pelo sistema operacional ou pelo próprio SGBD/sistema de arquivos (como o Linux mdadm, o Windows Storage Spaces ou o ZFS), usando a CPU do host para o cálculo de paridade --- opção mais barata, porém tipicamente com maior latência e uso de CPU sob carga de escrita intensa.

\subsection{Quadro Comparativo das Interfaces DAS}
\begin{table}[H]\centering\caption{Comparativo das interfaces de conexão para DAS discutidas no enunciado.}\label{tab:interfaces}
\small\begin{tabularx}{\textwidth}{@{}p{2cm}p{1.5cm}p{2cm}p{2cm}p{2cm}Y@{}}\toprule
\textbf{Interface} & \textbf{Geração/Ano} & \textbf{Velocidade típica} & \textbf{Topologia} & \textbf{Multi-iniciador / Dual-path} & \textbf{Uso típico em SBD}\\\midrule
SATA III & 2009 & 6 Gbps ($\approx$ 600 MB/s) & Ponto a ponto & Não & Armazenamento de baixo custo, servidores de entrada\\
eSATA & 2004 & 6 Gbps (idêntico ao SATA) & Ponto a ponto (externo) & Não & Discos externos legados (em declínio face ao USB)\\
SAS-4 (24G) & 2019 & 22,5 Gbps ($\approx$ 2.400 MB/s) & Ponto a ponto com expansores & Sim (dual-port) & Servidores de banco de dados e arrays RAID corporativos\\
FireWire 800 & 2002 & 0,8 Gbps (800 Mbps) & Barramento peer-to-peer, encadeável & N/A (peer-to-peer) & Legado --- praticamente descontinuado\\
USB 3.2 Gen2x2 & 2017 & 20 Gbps & Ponto a ponto (host-cêntrico) & Não & Discos externos, backup, DAS portátil\\
USB4 v2.0 & 2022 & Até 80 Gbps (120 Gbps assimétrico) & Ponto a ponto (host-cêntrico), túnel PCIe & Não & DAS externo de altíssimo desempenho\\
HBA & --- & Depende da interface física agregada & Adaptador PCIe host $\leftrightarrow$ storage & Depende do modo (JBOD/RAID) & Ponte entre o barramento do servidor e discos SAS/SATA\\\bottomrule
\end{tabularx}\end{table}

\section{RAID: Conceito, Paralelismo e Redundância}
RAID (originalmente Redundant Array of Inexpensive Disks, hoje mais comumente lido como Redundant Array of Independent Disks) é a técnica de combinar dois ou mais discos físicos independentes em um único volume lógico, gerenciado de forma transparente para o sistema operacional e para o SGBD. O conceito foi formalizado em 1988 por David Patterson, Garth Gibson e Randy Katz, na Universidade da Califórnia em Berkeley, motivados pela constatação de que um conjunto de discos pequenos e baratos, organizados de forma inteligente, poderia superar em desempenho e em custo-benefício um único disco grande e caro --- desde que se resolvesse o problema de que múltiplos discos, tomados individualmente, falham com mais frequência do que um único disco.

Cada nível de RAID representa um ponto de equilíbrio diferente entre dois objetivos que, em geral, competem entre si pelo mesmo recurso (o número de discos físicos disponíveis): paralelismo, voltado a desempenho, e redundância, voltada a disponibilidade e confiabilidade.

\subsection{Paralelismo (Desempenho)}
Paralelismo, no contexto de RAID, significa distribuir as operações de leitura e escrita entre vários discos, de modo que eles trabalhem simultaneamente sobre uma mesma requisição lógica (ou sobre requisições distintas concorrentes), em vez de um único disco atender a tudo sequencialmente. Isso é obtido por meio da técnica de striping (detalhada na Seção 6.1), que fragmenta os dados entre os discos do arranjo. O ganho de desempenho é, no caso ideal, proporcional ao número de discos: um arranjo com N discos pode, em tese, entregar até N vezes a taxa de transferência sequencial de um único disco, além de permitir que múltiplas operações de I/O pequenas e independentes --- típicas de cargas OLTP em um SBD --- sejam atendidas em paralelo por discos diferentes, reduzindo o tempo médio de resposta.

\subsection{Redundância (Disponibilidade e Confiabilidade)}
Redundância consiste em armazenar informação extra --- uma cópia completa dos dados (mirroring) ou informação de paridade calculada a partir deles --- de modo que a falha de um disco (ou, em alguns níveis, de até dois discos simultaneamente) não cause perda de dados nem interrupção do serviço. Disponibilidade refere-se à capacidade do arranjo de continuar respondendo a requisições mesmo em estado degradado (com um disco falho, antes da substituição); confiabilidade refere-se à probabilidade de que os dados permaneçam corretos e recuperáveis ao longo do tempo, mesmo diante de falhas de hardware.

Paralelismo puro (striping sem nenhuma redundância) reduz a confiabilidade do arranjo em relação a um único disco: como qualquer falha em qualquer um dos N discos torna o volume lógico inteiro inacessível, o MTTF (tempo médio até a falha) do arranjo é, em uma aproximação simples, o MTTF de um disco individual dividido por N --- ou seja, quanto mais discos participam do striping, maior a chance de que algum deles falhe em um dado intervalo de tempo. Os níveis de RAID que envolvem redundância existem precisamente para compensar esse efeito, reintroduzindo tolerância a falhas ao custo de capacidade útil (no caso do mirroring) ou de operações adicionais de escrita para manter a paridade atualizada (no caso dos níveis baseados em paridade). A escolha de um nível de RAID é, portanto, sempre uma decisão de engenharia que pondera paralelismo (desempenho) contra redundância (disponibilidade/confiabilidade) e contra o custo em capacidade de disco disponível.

\section{Organização de Dados em RAID: Striping, Mirroring e Paridade}
\subsection{Data Striping: Bit-Level e Block-Level}
Striping é a técnica de fragmentar um fluxo lógico de dados em pequenas unidades (chamadas de stripe units ou chunks) e distribuí-las ciclicamente (round-robin) entre os discos do arranjo, de modo que discos diferentes armazenem partes diferentes de um mesmo conjunto de dados e possam ser acessados em paralelo. O tamanho da unidade de striping determina duas variantes com comportamentos bem distintos:
\begin{itemize}[leftmargin=*]
    \item \textbf{Bit-level striping:} cada unidade de striping é um único bit (na prática, um pequeno agrupamento de bits), distribuído entre todos os discos do arranjo, que precisam girar de forma sincronizada (lockstep) --- qualquer acesso lógico, por menor que seja, envolve necessariamente todos os discos ao mesmo tempo. Essa técnica maximiza a taxa de transferência para transferências sequenciais grandes, mas é ineficiente para cargas de I/O pequeno e aleatório (típicas de OLTP), pois não permite que discos diferentes atendam requisições independentes em paralelo. É a base dos níveis RAID2 e RAID3.
    \item \textbf{Block-level striping:} cada unidade de striping é um bloco inteiro (tipicamente de alguns KB a alguns MB --- o "tamanho de stripe" ou chunk size, um parâmetro configurável), atribuído de forma independente a cada disco. Como cada bloco reside inteiramente em um único disco, requisições de I/O pequenas e não relacionadas podem ser atendidas por discos diferentes ao mesmo tempo, o que é essencial para o padrão de acesso aleatório típico de um SBD transacional. É a técnica usada nos níveis RAID0, RAID4, RAID5, RAID6 e nos arranjos aninhados RAID0+1/RAID10.
\end{itemize}

A escolha do tamanho do stripe é, na prática, um parâmetro de ajuste (tuning): stripes menores favorecem cargas com muitas operações pequenas e concorrentes (ex.: OLTP), pois espalham melhor as requisições entre os discos; stripes maiores favorecem cargas dominadas por leituras sequenciais de grandes volumes (ex.: varreduras completas de tabela em um Data Warehouse), pois reduzem o número de discos que precisam ser acionados para satisfazer uma única leitura sequencial longa.

\subsection{Mirroring e Shadowing}
Mirroring (espelhamento) consiste em manter uma cópia idêntica e simultânea dos dados em dois (ou mais) discos: toda escrita é replicada em ambos os discos do par, e toda leitura pode, em princípio, ser atendida por qualquer um dos dois --- o que permite, inclusive, balancear leituras entre as cópias para ganhar desempenho de leitura. É a técnica central do RAID1: garante tolerância à falha de um dos discos do par, às custas de um custo de capacidade de 50\% (para um espelho de dois discos, apenas metade do espaço físico total fica disponível como capacidade útil).

Shadowing ("espelhamento em sombra") é um termo historicamente empregado, sobretudo na literatura mais antiga de sistemas de banco de dados e em sistemas tolerantes a falhas como o Tandem NonStop, como sinônimo de mirroring: manter um disco "sombra" (shadow disk) que replica continuamente o conteúdo de um disco primário, de forma transparente para as aplicações. Na prática, os dois termos descrevem o mesmo mecanismo --- duplicação síncrona e contínua dos dados em um segundo disco físico ---, e é comum encontrá-los usados de forma intercambiável tanto na bibliografia da disciplina quanto na documentação de fabricantes; a diferença é, essencialmente, de origem histórica/vendor, não de técnica.

\subsection{Paridade e o Bloco de Paridade}
Paridade é uma forma de redundância mais eficiente em espaço do que o mirroring completo: em vez de duplicar integralmente os dados, calcula-se, a partir de N blocos de dados correspondentes em N-1 discos, um bloco de paridade adicional --- tipicamente por meio da operação lógica XOR (ou exclusive-or) bit a bit entre os blocos --- e armazena-se esse bloco de paridade em um disco separado (RAID3/RAID4) ou distribuído entre todos os discos do arranjo (RAID5/RAID6). Como a operação XOR é reversível (o valor de qualquer um dos blocos originais pode ser recalculado a partir dos demais blocos de dados e do bloco de paridade), a perda de um único disco pode ser totalmente reconstruída.

O custo dessa eficiência de espaço aparece nos escritas: sempre que um bloco de dados é alterado, o bloco de paridade correspondente também precisa ser recalculado e regravado, o que --- para escritas pequenas que não cobrem uma faixa (stripe) inteira --- exige um padrão de leitura e escrita adicional conhecido como read-modify-write (ler o bloco de dados antigo, ler o bloco de paridade antigo, calcular a nova paridade, gravar o novo bloco de dados e gravar a nova paridade). Esse é o chamado "write penalty" dos níveis baseados em paridade, discutido em detalhe na Seção 8 para cada nível específico.

\section{Métricas de Confiabilidade: MTTF, MTBF, MTTR e MTTDL}
Para comparar e dimensionar arranjos RAID de forma quantitativa, e não apenas qualitativa, a literatura de sistemas de armazenamento define um conjunto de métricas de confiabilidade. Estas são as quatro métricas exigidas pelo enunciado da tarefa.

\subsection{MTTF (Mean Time To Failure) e MTBF (Mean Time Between Failures)}
MTTF (tempo médio até a falha) é o tempo médio esperado de operação de um componente até sua primeira falha; é a métrica apropriada para itens que, uma vez falhos, são substituídos (não reparados), como é o caso típico de um disco individual em um servidor. MTBF (tempo médio entre falhas), formalmente, aplica-se a sistemas reparáveis e corresponde à soma do MTTF com o MTTR (tempo médio de reparo): MTBF = MTTF + MTTR. Na prática do mercado de armazenamento, no entanto, os dois termos são frequentemente usados como sinônimos em datasheets de fabricantes de HDD/SSD --- que costumam divulgar valores estatísticos da ordem de 1 a 2,5 milhões de horas --- já que o MTTR de um disco individual (o tempo até que ele seja fisicamente trocado) tende a ser muito pequeno frente ao seu MTTF, tornando a diferença entre as duas métricas pouco relevante nesse contexto específico. É importante frisar que esse número é estatístico (obtido a partir de testes acelerados e de grandes populações de discos), não uma garantia de vida útil para uma unidade específica.

\subsection{MTTR (Mean Time To Repair)}
MTTR (tempo médio de reparo) é o tempo médio necessário para detectar a falha de um disco, substituí-lo fisicamente (ou acionar um hot spare já instalado no arranjo) e concluir o processo de reconstrução (rebuild) dos dados no novo disco a partir dos discos sobreviventes e da informação de paridade ou espelho. O MTTR é especialmente relevante em RAID porque, durante toda a janela de reconstrução, o arranjo opera em estado degradado --- sem a redundância completa ---, e uma segunda falha nesse intervalo pode causar perda de dados. Discos modernos de alta capacidade (16 TB, 20 TB ou mais) tornam o MTTR um fator cada vez mais crítico: o tempo necessário para ler dezenas de terabytes dos discos sobreviventes e regravar o disco substituto pode facilmente ultrapassar 24 horas, ampliando significativamente a janela de vulnerabilidade do arranjo.

\subsection{MTTDL (Mean Time To Data Loss)}
MTTDL (tempo médio até a perda de dados) é a métrica que sintetiza, para um arranjo redundante, o tempo médio esperado até que a redundância seja de fato esgotada e ocorra perda efetiva de dados --- ou seja, até que uma segunda falha (independente da primeira) ocorra dentro da janela de vulnerabilidade aberta pela primeira falha e pelo seu MTTR. Para um RAID5 com N discos, uma aproximação comumente usada na literatura é:
\[ \text{MTTDL} \approx \frac{\text{MTTF}^2}{N \times (N-1) \times \text{MTTR}} \]
Essa fórmula evidencia três relações importantes para o dimensionamento de um SBD: (i) o MTTDL cresce com o quadrado do MTTF de cada disco individual, ou seja, discos mais confiáveis têm impacto desproporcionalmente positivo na confiabilidade do arranjo; (ii) o MTTDL cai à medida que o número de discos N aumenta, já que arranjos maiores têm mais oportunidades de uma segunda falha ocorrer; e (iii) o MTTDL cai linearmente com o aumento do MTTR, o que explica por que reduzir o tempo de reconstrução (por exemplo, com hot spares já instalados e prontos, reduzindo o tempo até o início do rebuild) é uma prática recomendada. É justamente essa sensibilidade ao MTTR --- agravada pelos discos de alta capacidade atuais --- que motiva o uso de RAID6 (tolerante a duas falhas simultâneas) em arranjos grandes, como discutido na Seção 8.7.

\section{Níveis Padrão de RAID (0 a 6, 0+1 e 10)}
Esta seção descreve, individualmente, cada um dos níveis padrão de RAID definidos na literatura clássica (0 a 6), além dos arranjos aninhados RAID0+1 e RAID10, encerrando com um quadro comparativo.

\subsection{RAID 0 --- Striping}
RAID0 aplica block-level striping puro, sem nenhuma redundância. Com N discos, a capacidade útil é N vezes a capacidade de um disco, e o desempenho de leitura e escrita é, no caso ideal, até N vezes o de um único disco --- o maior desempenho entre todos os níveis. Em compensação, é o único nível menos confiável do que um disco isolado: a falha de qualquer disco do arranjo torna todo o volume lógico inacessível, pois não há informação redundante para reconstruí-lo. É recomendado apenas para dados não críticos, facilmente reproduzíveis a partir de outra fonte, ou protegidos por redundância em outra camada --- por exemplo, áreas de trabalho temporárias (tempdb/scratch space) de um SGBD.

\subsection{RAID 1 --- Mirroring}
RAID1 aplica espelhamento total entre (tipicamente) dois discos: cada byte gravado em um disco é replicado integralmente no outro. A capacidade útil equivale à de um único disco (50\% de overhead), e o arranjo tolera a falha de um dos dois discos sem qualquer perda de dados. O desempenho de escrita é próximo ao de um único disco (ambos precisam ser gravados), mas o desempenho de leitura pode ser superior, já que requisições de leitura podem ser distribuídas entre as duas cópias. A reconstrução após uma falha é simples e relativamente rápida: basta copiar integralmente o disco sobrevivente para o disco substituto. É o nível recomendado para volumes pequenos e críticos, como os logs de transação (redo/undo log) e volumes de sistema de um SGBD.

\subsection{RAID 2 --- Bit-Level Striping com Código Hamming}
RAID2 aplica bit-level striping combinado a um código de correção de erros Hamming, calculado sobre os bits correspondentes em cada disco e armazenado em discos dedicados de verificação (checksum), permitindo não apenas detectar, mas corrigir determinados padrões de erro sem a necessidade de um disco de paridade separado como nos níveis seguintes. Exige um número relativamente grande de discos sincronizados (por exemplo, uma configuração clássica usa 10 discos de dados e 4 discos de Hamming) girando em lockstep. Na prática, o RAID2 é hoje um nível essencialmente histórico/teórico: como os discos modernos (HDD e SSD) já incorporam sua própria correção de erro em nível de setor (ECC embutido no firmware do disco), a camada adicional de Hamming do RAID2 tornou-se redundante, e não há implementações comerciais relevantes desse nível atualmente.

\subsection{RAID 3 --- Byte/Bit-Level Striping com Paridade Dedicada}
RAID3 usa striping em nível de byte (ou bit) com um único disco dedicado exclusivamente à paridade. Como todos os discos precisam operar em lockstep para atender a uma única requisição (cada I/O lógico é fragmentado entre todos os discos), o RAID3 entrega taxas de transferência sequenciais muito altas --- todos os discos trabalham juntos em cada operação ---, mas desempenho de I/O aleatório baixo, já que não é possível atender duas requisições pequenas e independentes em paralelo (todos os discos, incluindo o de paridade, estão sempre ocupados com a mesma operação). É indicado para cargas dominadas por transferências sequenciais grandes de um único fluxo --- por exemplo, ingestão em lote (bulk load) ou aplicações de streaming/imagens ---, mas pouco adequado ao padrão de acesso concorrente e aleatório típico de OLTP.

\subsection{RAID 4 --- Block-Level Striping com Paridade Dedicada}
RAID4 aplica a mesma ideia de um disco de paridade dedicado do RAID3, mas com striping em nível de bloco em vez de byte/bit. Isso elimina a exigência de sincronismo entre discos: como cada bloco reside inteiramente em um disco, requisições pequenas e independentes podem ser atendidas por discos diferentes em paralelo, o que melhora bastante o desempenho de leitura aleatória em relação ao RAID3. O grande problema do RAID4, porém, é que toda escrita --- em qualquer disco de dados do arranjo --- precisa também atualizar o único disco de paridade, que se torna um gargalo de serialização sob carga de escrita concorrente (o "parity disk bottleneck"). Por essa limitação, o RAID4 é hoje pouco utilizado, tendo sido amplamente substituído pelo RAID5, que resolve exatamente esse gargalo.

\subsection{RAID 5 --- Block-Level Striping com Paridade Distribuída}
RAID5 resolve o gargalo do RAID4 distribuindo (fazendo rotacionar) os blocos de paridade entre todos os discos do arranjo, em vez de concentrá-los em um único disco dedicado --- cada disco armazena, para diferentes faixas (stripes), tanto blocos de dados quanto blocos de paridade. Isso distribui também a carga de escrita de paridade entre todos os discos, eliminando o ponto único de serialização do RAID4. Exige no mínimo 3 discos; a capacidade útil é (N-1) vezes a capacidade de um disco (o equivalente de um disco é "gasto" em paridade, distribuído entre todos). RAID5 tolera a falha de um único disco. Seu principal custo é o já mencionado write penalty: uma escrita pequena que não cobre uma faixa inteira ainda exige um ciclo de read-modify-write (4 operações de I/O: ler dado antigo, ler paridade antiga, gravar novo dado, gravar nova paridade). Ainda assim, é hoje o nível de RAID mais utilizado em SBD de uso geral, por equilibrar bem capacidade, desempenho e tolerância a falha --- com uma ressalva importante: em discos de alta capacidade, o tempo de reconstrução após uma falha pode ser longo (ver Seção 7.2 e 7.3), elevando o risco de uma segunda falha durante o rebuild, o que tem levado à migração para RAID6 em arranjos grandes.

\subsection{RAID 6 --- Dupla Paridade Distribuída}
RAID6 estende o RAID5 com um segundo bloco de paridade, independente do primeiro (tipicamente calculado com um esquema do tipo Reed-Solomon, formando um par P+Q), também distribuído entre todos os discos. Isso permite tolerar a falha simultânea de dois discos quaisquer do arranjo. Exige no mínimo 4 discos; a capacidade útil é (N-2) vezes a capacidade de um disco. O write penalty é ainda maior do que o do RAID5 (até 6 operações de I/O por escrita pequena: ler dado antigo, ler as duas paridades antigas, gravar o novo dado e as duas novas paridades), mas o ganho em MTTDL é substancial, especialmente em arranjos com muitos discos de alta capacidade e MTTR longo --- cenário em que uma segunda falha durante a reconstrução deixou de ser um risco desprezível. Por isso, o RAID6 é hoje amplamente recomendado para arranjos grandes e de alta capacidade em SBD de produção, mesmo com o custo adicional em desempenho de escrita e em capacidade útil.

\subsection{RAID 0+1 e RAID 10 --- Arranjos Aninhados de Striping e Mirroring}
Ambos combinam striping (RAID0) e mirroring (RAID1) para obter, simultaneamente, alto desempenho e redundância sem o write penalty dos níveis baseados em paridade --- ao custo de 50\% da capacidade total, como no RAID1 puro. A diferença entre os dois está na ordem em que as técnicas são combinadas, e essa ordem tem consequências importantes para a confiabilidade:
\begin{itemize}[leftmargin=*]
    \item \textbf{RAID 0+1 (mirror de stripes):} primeiro cria-se dois conjuntos RAID0 (stripes) e, depois, espelha-se um conjunto inteiro contra o outro. Se um único disco falha em um dos conjuntos, todo aquele conjunto stripe é considerado comprometido (pois o RAID0 não tem redundância própria), restando como única proteção o conjunto espelhado íntegro; uma segunda falha em qualquer disco do conjunto sobrevivente --- mesmo em um disco diferente do primeiro --- causa perda total dos dados. Além disso, a reconstrução exige recopiar o conjunto inteiro (todo o stripe), não apenas o disco individual falho.
    \item \textbf{RAID 10 (ou RAID 1+0, stripe de espelhos):} primeiro cria-se pares espelhados (RAID1) e, depois, faz-se striping entre esses pares. Uma falha de disco compromete apenas o par ao qual ele pertence; os demais pares continuam plenamente espelhados. Uma segunda falha só causa perda de dados se atingir justamente o disco parceiro do primeiro disco falho dentro do mesmo par --- em qualquer outro par, o arranjo continua tolerando novas falhas. A reconstrução, além disso, é muito mais rápida, pois exige copiar apenas o disco individual falho a partir de seu par, não o conjunto inteiro. Por essas razões, o RAID10 é considerado, na prática, superior ao RAID0+1, sendo o padrão recomendado para bancos de dados relacionais com alta taxa de transação (OLTP) em que tanto desempenho quanto confiabilidade são críticos.
\end{itemize}

\subsection{Quadro Comparativo dos Níveis Padrão de RAID}
\begin{longtable}{@{}p{1.5cm}p{1.2cm}p{2.5cm}p{1.5cm}p{1.8cm}p{1.8cm}Y@{}}\toprule
\textbf{Nível} & \textbf{Nº mín. discos} & \textbf{Técnica} & \textbf{Discos tolerados} & \textbf{Capacidade útil} & \textbf{Desempenho escrita} & \textbf{Uso recomendado em SBD}\\\midrule\endhead
RAID 0 & 2 & Striping em bloco & 0 & $N \times \text{disco}$ & Muito alto & Dados não críticos/temporários (tempdb)\\
RAID 1 & 2 & Espelhamento (mirroring) & 1 (de 2) & $1 \times \text{disco}$ & Moderado (2 gravações) & Logs de transação, volumes de sistema\\
RAID 2 & Muitos (ex.: 10+4) & Striping em bit + Hamming & 1 (teórico) & $(N-\text{check}) \times \text{disco}$ & Alto (teórico) & Obsoleto --- sem uso comercial atual\\
RAID 3 & 3 & Striping byte + paridade dedicada & 1 & $(N-1) \times \text{disco}$ & Baixo (gargalo na paridade) & Transferências sequenciais grandes (streaming, bulk load)\\
RAID 4 & 3 & Striping em bloco + paridade dedicada & 1 & $(N-1) \times \text{disco}$ & Baixo (gargalo na paridade) & Pouco usado --- superado pelo RAID5\\
RAID 5 & 3 & Striping em bloco + paridade distribuída & 1 & $(N-1) \times \text{disco}$ & Moderado (write penalty) & Uso geral em SBD; equilíbrio custo $\times$ desempenho\\
RAID 6 & 4 & Striping em bloco + dupla paridade distribuída & 2 & $(N-2) \times \text{disco}$ & Moderado-baixo & Arranjos grandes/discos de alta capacidade\\
RAID 0+1 & 4 & Espelho de dois stripes & 1 (risco elevado na 2ª falha) & 50\% & Alto & Pouco recomendado --- RAID10 é preferível\\
RAID 10 & 4 & Stripe de espelhos & Múltiplos (pares distintos) & 50\% & Alto & OLTP de alta transação --- padrão para dados críticos\\\bottomrule
\end{longtable}
Tabela 2 --- Comparativo dos níveis padrão de RAID (0 a 6, 0+1 e 10).

\section{Níveis Não Padrão de RAID}
Além dos níveis padrão descritos na Seção 8 --- que são conceitos abertos, definidos academicamente e não amarrados a um único fabricante ---, o mercado de armazenamento desenvolveu, ao longo dos anos, diversas implementações proprietárias e arranjos aninhados que estendem ou combinam esses conceitos de formas específicas. O enunciado da tarefa pede o estudo de RAID 1.5, RAID 7, RAID-DP, RAID-S e Matrix RAID, discutidos a seguir, além de uma menção aos arranjos aninhados de maior escala (RAID 50/60/100).

\subsection{RAID 1.5}
RAID 1.5 é uma tecnologia proprietária da HighPoint Technologies (implementada em suas controladoras RocketRAID), que combina, em um arranjo de apenas dois discos, características de RAID0 e RAID1 simultaneamente: as escritas são espelhadas integralmente entre os dois discos (como no RAID1, garantindo redundância total), enquanto as leituras são divididas/distribuídas entre os dois discos (como no RAID0, ganhando desempenho de leitura paralela). O objetivo é obter, com apenas dois discos, uma leitura mais rápida do que um RAID1 convencional ofereceria, sem pagar o custo de quatro discos que um RAID10 exigiria. Por ser uma tecnologia específica de um fabricante, não é um padrão aberto nem amplamente adotado fora do ecossistema de controladoras daquela empresa.

\subsection{RAID 7}
RAID 7 é uma marca registrada da (hoje extinta) Storage Computer Corporation, e não faz parte da série acadêmica/aberta de níveis de RAID --- apesar do nome sugerir uma continuação natural do RAID6, trata-se de uma arquitetura proprietária específica dessa empresa. Sua proposta central era eliminar os gargalos de sincronismo dos primeiros níveis (RAID3/RAID4) por meio de operação assíncrona: cada disco do arranjo possuía sua própria cache independente, um microprocessador embarcado com sistema operacional de tempo real coordenava as operações do arranjo (em vez de depender apenas do host), e um barramento de alta velocidade dedicado, separado do barramento do host, interligava os discos e o controlador. Por não ser um padrão aberto --- e por a empresa que o registrou não existir mais ---, o RAID 7 é hoje essencialmente uma curiosidade histórica, não uma opção disponível no mercado atual.

\subsection{RAID-DP (Double Parity)}
RAID-DP é a implementação proprietária de dupla paridade da NetApp, integrada ao seu sistema de arquivos WAFL (Write Anywhere File Layout) e ao sistema operacional de armazenamento Data ONTAP. Conceitualmente, é semelhante ao RAID6 acadêmico --- tolera a falha simultânea de até dois discos ---, mas usa duas paridades com papéis diferentes: uma paridade horizontal (semelhante à de um RAID4/5 tradicional, calculada sobre uma faixa/stripe) e uma segunda paridade diagonal, calculada sobre uma diagonal de blocos que atravessa múltiplas faixas. A principal vantagem alegada pela NetApp é que essa combinação, aliada ao design do WAFL --- que sempre escreve em blocos livres em vez de sobrescrever no local (evitando o clássico read-modify-write) ---, reduz a penalidade de escrita normalmente associada a esquemas de dupla paridade, tornando o RAID-DP mais eficiente em operações de escrita do que uma implementação convencional de RAID6.

\subsection{RAID-S}
RAID-S é um nome historicamente usado pela EMC para uma implementação de paridade proprietária em seus arrays de armazenamento Symmetrix, funcionalmente próxima de um RAID3/RAID5 tradicional (striping com paridade), porém otimizada para a arquitetura específica desses arrays --- em particular, para a cache espelhada e o barramento interno de alta velocidade da EMC, usados para amenizar o custo do read-modify-write da paridade. Como no caso do RAID-DP e do RAID 1.5, trata-se de uma tecnologia associada a uma linha de produtos de um fabricante específico, e não de um padrão aberto definido independentemente de implementação, como os níveis 0 a 6 descritos na Seção 8.

\subsection{Matrix RAID (Intel)}
Matrix RAID --- comercializado pela Intel como parte de sua Intel Matrix Storage Technology, hoje parte do Intel Rapid Storage Technology (RST) --- não é, tecnicamente, um novo nível de RAID, mas sim um recurso que permite criar dois (ou mais) volumes RAID de níveis diferentes usando o mesmo conjunto físico de discos. Um exemplo típico usa apenas dois discos físicos: cada disco é particionado internamente em duas regiões, e o controlador Intel mapeia uma região de cada disco como um pequeno volume RAID0 (para o sistema operacional e aplicações, priorizando desempenho) e a outra região como um volume RAID1 (para dados críticos do usuário, priorizando redundância) --- obtendo parte dos benefícios de desempenho e de proteção sem exigir quatro discos físicos separados, como um RAID10 tradicional exigiria. É um recurso voltado sobretudo a desktops e estações de trabalho equipados com chipsets Intel, mais do que a servidores de banco de dados de produção.

\subsection{Outros Arranjos Aninhados: RAID 50, RAID 60 e RAID 100}
Além dos casos proprietários acima, é comum na indústria (embora sem um único órgão de padronização formal) combinar níveis padrão em dois estágios para arranjos de grande porte: RAID50 faz striping (RAID0) entre múltiplos subarranjos RAID5; RAID60 faz striping entre múltiplos subarranjos RAID6; e RAID100 (ou RAID1+0+0) faz striping entre múltiplos subarranjos RAID10. Esses arranjos aninhados são usados sobretudo em SANs e subsistemas de armazenamento corporativos de grande escala, quando um único subarranjo (por exemplo, um único grupo RAID6) atingiria um número de discos impraticável para os tempos de rebuild envolvidos: dividir os discos em múltiplos subgrupos menores, unidos por striping, reduz o tempo de reconstrução de cada falha individual e melhora o desempenho agregado, mantendo a tolerância a falha de cada subgrupo isoladamente.

\section{Recomendações de Uso e Aplicação em SBD}
Combinando os conceitos discutidos nas seções anteriores, é possível formular recomendações práticas de arquitetura de armazenamento para diferentes componentes típicos de um Sistema de Banco de Dados:
\begin{itemize}[leftmargin=*]
    \item \textbf{Logs de transação (redo/undo log):} beneficiam-se de RAID1 ou RAID10 sobre discos SAS (ou NVMe, fora do escopo deste trabalho): a prioridade é baixa latência de escrita e redundância, sem o write penalty de paridade, já que o log é escrito de forma predominantemente sequencial e sua perda compromete a durabilidade (propriedade D do ACID) de todas as transações desde o último checkpoint.
    \item \textbf{Arquivos de dados/tablespaces de OLTP:} RAID10 é, em geral, a recomendação padrão quando o orçamento permite, por combinar o melhor desempenho de I/O aleatório (sem write penalty) com boa tolerância a falhas; RAID5 ou RAID6 são alternativas mais econômicas em capacidade quando o padrão de acesso é menos intensivo em escritas aleatórias pequenas.
    \item \textbf{Data Warehouse / cargas analíticas (OLAP):} como o padrão de acesso é predominantemente sequencial (varreduras de grandes volumes de dados) e o custo por gigabyte tende a ser prioridade frente à latência de escrita individual, RAID5 ou RAID6 costumam oferecer boa relação entre capacidade útil e confiabilidade; RAID6 é preferível quando os discos são de alta capacidade, dado o risco de uma segunda falha durante o rebuild.
    \item \textbf{Backup e arquivamento de longo prazo:} RAID6 é recomendado por sua tolerância a duas falhas simultâneas, essencial quando os dados precisam ser retidos por longos períodos e o tempo entre verificações/rebuilds pode ser maior.
    \item \textbf{Áreas temporárias/scratch (tempdb):} RAID0 é aceitável quando os dados são inteiramente reprodutíveis ou descartáveis em caso de falha, e o objetivo é exclusivamente desempenho.
\end{itemize}

Quanto à interface e à controladora, discos SAS com uma controladora RAID de hardware dotada de cache protegida por bateria (BBU) ou flash (FBWC) são a escolha típica para servidores de produção, por combinar dual-porting, fila de comandos profunda e certificação para operação contínua; SATA permanece adequado para cenários de custo mais sensível ou de menor criticidade; e USB4 (sobretudo a partir da Versão 2.0, com até 80 Gbps) vem se tornando uma opção cada vez mais viável para DAS externo de alto desempenho, como estações de backup ou ambientes de desenvolvimento/teste. eSATA e FireWire, por sua vez, são hoje interfaces predominantemente legadas, raramente escolhidas para novos projetos de SBD.

\section{Conclusão}
O estudo das arquiteturas de armazenamento físico para SBD evidenciou que a escolha de uma interface de conexão (SATA, eSATA, SAS, FireWire, USB ou o HBA que as agrega ao host) e de um nível de RAID não são decisões puramente técnicas de baixo nível, mas escolhas de arquitetura com impacto direto sobre desempenho, disponibilidade e confiabilidade de um banco de dados --- e, em última instância, sobre a experiência dos usuários e a integridade dos dados sob a responsabilidade do SGBD.

Os conceitos de paralelismo e redundância mostraram-se, ao longo do trabalho, dois lados de uma mesma moeda: técnicas como o striping em nível de bloco maximizam desempenho, enquanto mirroring e paridade sacrificam capacidade ou desempenho de escrita em troca de tolerância a falhas --- e métricas como MTTF, MTBF, MTTR e, sobretudo, MTTDL permitem quantificar esse equilíbrio de forma objetiva, orientando decisões como a preferência crescente por RAID6 (em vez de RAID5) em arranjos grandes de discos de alta capacidade, ou por RAID10 (em vez de RAID0+1) sempre que o orçamento de capacidade permite. Os níveis não padrão --- RAID 1.5, RAID 7, RAID-DP, RAID-S e Matrix RAID --- reforçam que, na prática comercial, fabricantes frequentemente adaptam ou estendem os conceitos acadêmicos de RAID às particularidades de suas próprias arquiteturas de hardware e software, o que exige cautela ao comparar especificações entre fornecedores diferentes.

Para um profissional de Banco de Dados, compreender essas arquiteturas --- desde o cabo físico até a métrica estatística de confiabilidade --- é essencial para dimensionar corretamente a infraestrutura de armazenamento de um sistema, equilibrando custo, desempenho e, acima de tudo, a segurança dos dados armazenados.

\section{Post-Mortem --- Uso de Assistente de IA Generativa}
Conforme autorizado e exigido pelo enunciado da tarefa, esta seção documenta de forma auditável o uso de assistente de IA generativa (Claude, da Anthropic) na elaboração deste trabalho: a metodologia de uso, o log cronológico dos prompts-chave, a matriz de responsabilidades e autoria dos integrantes do grupo, a declaração individual de autoria, os erros e riscos identificados no conteúdo gerado pela IA (com a respectiva correção) e um checklist final de conformidade com cada exigência do enunciado.

\subsection{Metodologia de Uso da IA}
O grupo adotou um fluxo de trabalho "humano no comando" (human-in-the-loop), em quatro etapas:
\begin{enumerate}
    \item Delimitação do escopo e do formulário --- o grupo forneceu à IA o enunciado completo da tarefa sobre DAS e RAID, além dos arquivos de um trabalho anterior da mesma disciplina, para que a IA seguisse a mesma estrutura, formatação e nível de profundidade já validados anteriormente com o professor.
    \item Geração assistida com pesquisa ativa --- para dados técnicos sujeitos a mudança --- como as velocidades atuais de SAS-4/24G e da especificação USB4 Versão 2.0 ---, a IA foi instruída a buscar informações atualizadas na web em vez de responder apenas com base em conhecimento de treinamento, reduzindo o risco de valores desatualizados sobre gerações/velocidades de interfaces que evoluem com frequência.
    \item Revisão crítica humana --- cada integrante deve revisar as seções sob sua responsabilidade (ver Seção 12.3), verificando a coerência técnica frente ao material da disciplina e frente à documentação técnica dos fabricantes citados.
    \item Registro e rastreabilidade --- todas as decisões, prompts e eventuais correções foram/devem ser documentados nas subseções seguintes, permitindo que o professor e os colegas auditem exatamente o que foi gerado pela IA e o que foi validado/corrigido pelo grupo.
\end{enumerate}

\subsection{Log Cronológico dos Prompts-Chave Utilizados}
O quadro abaixo relaciona, em ordem cronológica, os prompts principais utilizados na conversa com a IA para a produção deste documento e dos slides. O grupo deve complementar este log com quaisquer prompts adicionais usados na revisão final, no preenchimento da folha de rosto e na preparação da apresentação.

\begin{longtable}{@{}p{.7cm}p{7.3cm}p{6.4cm}@{}}\toprule
\textbf{Nº} & \textbf{Prompt (resumo)} & \textbf{Objetivo}\\\midrule\endhead
1 & Envio do enunciado da tarefa sobre armazenamento físico DAS e RAID. & Estabelecer o escopo técnico completo exigido: interfaces DAS e níveis de RAID padrão/não padrão.\\
2 & Envio dos arquivos de um trabalho anterior da disciplina como referência de formatação. & Garantir consistência de estrutura, folha de rosto, seções e Post-Mortem com um padrão já validado pelo professor.\\
3 & Pedido para gerar um novo relatório (Word) e uma nova apresentação (PowerPoint) sobre DAS/RAID. & Produzir os dois artefatos finais exigidos pela tarefa em formato editável, mantendo o padrão visual da disciplina.\\
4 & Pesquisa na web sobre a velocidade atual do SAS-4 (24G) e da especificação USB4 Versão 2.0. & Fundamentar a Seção 4 com dados técnicos verificáveis e atuais, em vez de estimativas de memória do modelo.\\
5 & Geração do relatório completo (Word/PDF) e dos slides (PowerPoint), com folha de rosto, seções técnicas, tabelas comparativas e Post-Mortem inicial. & Entregar os dois produtos obrigatórios da tarefa já estruturados conforme o enunciado.\\\bottomrule
\end{longtable}
Tabela 3 --- Log de prompts-chave (a ser complementado pelo grupo com os prompts de sua própria revisão).

\subsection{Matriz de Responsabilidades e Autoria}
A tabela a seguir é um template de matriz de responsabilidades a ser preenchido pelo grupo, atribuindo a cada integrante a responsabilidade principal por seções específicas do trabalho.

\begin{longtable}{@{}p{4.2cm}p{5cm}p{5.2cm}@{}}\toprule
\textbf{Integrante} & \textbf{Seções sob sua responsabilidade} & \textbf{Decisão a justificar}\\\midrule\endhead
[Nome 1] --- DRE [....] & Seções 1, 2 e 3 (Resumo, Introdução, DAS/NAS/SAN) & Por que o DAS foi contextualizado frente a NAS e SAN antes de detalhar as interfaces.\\
[Nome 2] --- DRE [....] & Seção 4 (Interfaces DAS) e Tabela 1 & Por que cada interface foi classificada como indicada, ou não, para SBD de produção.\\
[Nome 3] --- DRE [....] & Seções 5, 6 e 7 (Paralelismo/Redundância; Striping/Mirroring/Paridade; Métricas) & Por que a fórmula de MTTDL apresentada foi escolhida e validada frente à bibliografia.\\
[Nome 4] --- DRE [....] & Seções 8 e 9 (Níveis Padrão e Não Padrão de RAID) e Tabela 2 & Por que cada nível não padrão foi classificado como proprietário/não aberto.\\
[Nome 5] --- DRE [....] & Seções 10, 11 e 12 (Recomendações; Conclusão; Post-Mortem) & Por que cada recomendação de uso foi associada ao componente de SBD correspondente.\\\bottomrule
\end{longtable}
Tabela 4 --- Matriz de responsabilidades e autoria (template a ser preenchido).

\subsection{Declaração Individual de Autoria}
Cada integrante deve preencher e confirmar a declaração abaixo, confirmando que revisou o conteúdo pelo qual é responsável.

\begin{longtable}{@{}p{5cm}p{4cm}p{5cm}@{}}\toprule
\textbf{Integrante} & \textbf{Declaro que revisei minha(s) seção(ões)} & \textbf{Data}\\\midrule\endhead
[Nome 1] & [ ] Sim & [__/__/2026]\\
[Nome 2] & [ ] Sim & [__/__/2026]\\
[Nome 3] & [ ] Sim & [__/__/2026]\\
[Nome 4] & [ ] Sim & [__/__/2026]\\
[Nome 5] & [ ] Sim & [__/__/2026]\\\bottomrule
\end{longtable}
Tabela 5 --- Declaração individual de autoria e ciência do conteúdo entregue.

\subsection{Erros e Riscos Identificados no Conteúdo Gerado pela IA e Correções Aplicadas}
\begin{longtable}{@{}p{.7cm}p{4.5cm}p{4.5cm}p{4.8cm}@{}}\toprule
\textbf{Nº} & \textbf{Risco / erro identificado} & \textbf{Como foi identificado} & \textbf{Correção aplicada}\\\midrule\endhead
1 & Velocidades de interfaces geradas "de memória" tendem a ficar desatualizadas. & Verificação ativa: a IA foi instruída a pesquisar fontes atuais. & Seção 4 fundamentada com valores de fontes técnicas recentes.\\
2 & Risco de confundir MTTF e MTBF como sinônimos absolutos. & Revisão crítica da Seção 7.1 frente à definição formal. & Texto explicita a diferença formal antes de mencionar uso de mercado.\\
3 & Níveis não padrão de RAID apresentados como se fossem abertos/uniformes. & Revisão comparando descrições com múltiplas fontes independentes. & Subseções explicitam fabricante/proprietário e ressalta que não é padrão aberto.\\
4 & Referências às seções exatas dos livros-texto não podem ser verificadas literalmente pela IA. & Reconhecimento explícito da limitação sobre direitos autorais. & Grupo deve confirmar correspondência exata na edição física/PDF.\\
5 & [Espaço para o grupo] & [Como o grupo identificou] & [Correção efetivamente aplicada]\\\bottomrule
\end{longtable}
Tabela 6 --- Riscos e limitações identificados no conteúdo gerado pela IA e correções aplicadas.

\subsection{Checklist de Conformidade com o Enunciado}
\begin{longtable}{@{}p{6.5cm}p{5cm}p{3cm}@{}}\toprule
\textbf{Exigência do enunciado} & \textbf{Onde é atendida} & \textbf{Status}\\\midrule\endhead
Interfaces DAS: SATA, eSATA, SAS, FireWire, USB (3 e 4), HBA & Seção 4 (4.1 a 4.6) e Tabela 1 & Atendido\\
Paralelismo (performance) e redundância (disponibilidade/confiabilidade) & Seção 5 (5.1 e 5.2) & Atendido\\
Data striping: bit-level e block-level striping & Seção 6.1 & Atendido\\
Mirroring e shadowing & Seção 6.2 & Atendido\\
MTTR, MTTDL, MTTF e MTBF & Seção 7 (7.1 a 7.3) & Atendido\\
Níveis de RAID padrão, bloco de paridade etc. & Seção 8 e Tabela 2; paridade em 6.3 & Atendido\\
RAID não padrão (1.5, 7, DP, S, Matrix etc.) & Seção 9 (9.1 a 9.6) & Atendido\\
Comparação de vantagens/desvantagens/recomendações & Seções 8, 9 e 10 & Atendido\\
Folha de rosto com nomes completos e DRE & Capa do documento & Pendente --- preencher\\
Log dos prompts-chave utilizados & Seção 12.2 e Tabela 3 & Atendido\\
Demonstração de autoria: quem fez o quê, quem decidiu o quê & Seções 12.3 e 12.4 & Pendente --- preencher\\
Identificação de erros da IA e correções aplicadas & Seção 12.5 e Tabela 6 & Atendido\\
Seção intitulada "Post-Mortem" & Seção 12 & Atendido\\
PDF postado na tarefa & Este documento, exportado em .pdf & Ação do grupo no AVA\\
Slides da apresentação postados & Arquivo .pptx elaborado em paralelo & Ação do grupo no AVA\\\bottomrule
\end{longtable}
Tabela 7 --- Checklist final de conformidade com o enunciado.

Este relatório Post-Mortem evidencia que o uso de IA generativa foi tratado de forma crítica e rastreável pelo grupo --- como ferramenta de produtividade e de primeira pesquisa, sempre sujeita à verificação humana e à divisão explícita de responsabilidades entre os integrantes.

\section{Referências Bibliográficas}
\begin{thebibliography}{99}
\bibitem{silberschatz} SILBERSCHATZ, A.; KORTH, H. F.; SUDARSHAN, S. \textit{Sistema de Banco de Dados}. 7ª ed. Rio de Janeiro: Elsevier, 2020. Seção 12.5 --- RAID.
\bibitem{elmasri} ELMASRI, R.; NAVATHE, S. B. \textit{Sistemas de Banco de Dados}. 7ª ed. São Paulo: Pearson, 2019. Seção 16.10 --- RAID.
\bibitem{garciamolina} GARCIA-MOLINA, H.; ULLMAN, J. D.; WIDOM, J. \textit{Database Systems: The Complete Book}. 2ª ed. Pearson, 2009. Seções 13.3 e 13.4 --- Disk Failures / RAID.
\bibitem{patterson} PATTERSON, D. A.; GIBSON, G.; KATZ, R. H. \textit{A Case for Redundant Arrays of Inexpensive Disks (RAID)}. Proceedings of the 1988 ACM SIGMOD International Conference on Management of Data.
\bibitem{snia} SNIA --- Storage Networking Industry Association. \textit{Common RAID Disk Data Format (DDF) Specification}.
\bibitem{scsi} SCSI Trade Association (STA). \textit{24G SAS Technology --- SAS-4 Overview}.
\bibitem{usbif} USB Implementers Forum (USB-IF). \textit{USB4® Version 2.0 Specification Announcement}. Disponível em: \url{https://www.usb.org/}. Acesso em: set. 2026.
\bibitem{netapp} NetApp, Inc. \textit{RAID-DP: NetApp Implementation of Double-Parity RAID for Data Protection}. Technical Report.
\bibitem{grupo} [Grupo]: incluir aqui quaisquer outras fontes efetivamente consultadas.
\end{thebibliography}

\end{document}